%% LyX 2.3.6.1 created this file.  For more info, see http://www.lyx.org/.
%% Do not edit unless you really know what you are doing.
\documentclass[english]{article}
\usepackage[latin9]{inputenc}
\usepackage[a4paper]{geometry}
\geometry{verbose}
\usepackage{color}
\usepackage{babel}
\usepackage{amsmath}
\usepackage{amssymb}
\usepackage{graphicx}
\usepackage[unicode=true,
 bookmarks=false,
 breaklinks=false,pdfborder={0 0 1},backref=section,colorlinks=true]
 {hyperref}
\hypersetup{
 linkcolor=blue,citecolor=blue,urlcolor=blue}

\makeatletter

%%%%%%%%%%%%%%%%%%%%%%%%%%%%%% LyX specific LaTeX commands.
%% Because html converters don't know tabularnewline
\providecommand{\tabularnewline}{\\}

%%%%%%%%%%%%%%%%%%%%%%%%%%%%%% User specified LaTeX commands.


% -------------------------------------------------------
% Packages
% -------------------------------------------------------
\usepackage{times}
\usepackage{amsfonts}\usepackage{bm}
\usepackage{bbm}
\usepackage[table]{xcolor}
\usepackage{multirow}
\usepackage{authblk}
\usepackage{caption}
%\usepackage{setspace}
%\singlespacing
% Double spacing for submission draft


% -------------------------------------------------------
% Convenience macros
% -------------------------------------------------------
\newcommand{\K}{\mathbf{K}}
\newcommand{\X}{\mathbf{X}}
\newcommand{\Hmat}{\mathbf{H}}
\newcommand{\Ps}{\mathbf{P}_{\mathrm{s}}}
\newcommand{\Pb}{\mathbf{P}_{\mathrm{b}}}
\newcommand{\Xs}{\mathbf{X}_{\mathrm{s}}}
\newcommand{\Xb}{\mathbf{X}_{\mathrm{b}}}
\newcommand{\Gs}{\boldsymbol{\Gamma}_{\mathrm{s}}}
\newcommand{\Gb}{\boldsymbol{\Gamma}_{\mathrm{b}}}
\newcommand{\lp}{\lambda_{+}}
\newcommand{\bstruct}{\hat{\beta}_{\mathrm{s}}}
\newcommand{\bbulk}{\hat{\beta}_{\mathrm{b}}}
\newcommand{\Lopt}{L_{\mathrm{opt}}}
\newcommand{\dtopt}{\Delta\tau_{\mathrm{opt}}}
\newcommand{\Fnorm}[1]{\left\|#1\right\|_{F}^{2}}
\newcommand{\E}{\mathbb{E}}
\newcommand{\tr}{\mathrm{Tr}}

% Roman section numbering as in Nature Physics
\renewcommand{\thesection}{\Roman{section}}
\renewcommand{\thesubsection}{\Alph{subsection}}

\author[1]{Rafael Morgado}
\author[2]{Ronne Toledo}
\author[1]{Eberth Correa}
\author[3]{Fernando A.\ Oliveira}
\affil[1]{Applied Natural Science Department, Faculty of Engineering Science and Technology, University of Bras�lia, 72444-240, Bras�lia, Federal District, Brazil}
\affil[2]{Aerospace Department, Faculty of Engineering Science and Technology, University of Bras\'ilia, 72444-240, Bras\'ilia, Federal District, Brazil}
\affil[3]{International Center of Physics, Institute of Physics, University of Bras\'ilia, 70910-900, Bras\'ilia, Federal District, Brazil}

\makeatother

\begin{document}
\title{\noindent Beyond Price Dynamics: Universal Exponents, Phase Transitions,
and a New Spectral Framework in Operator Space}
\maketitle
\begin{abstract}
\noindent We propose that the temporal covariance operator $\K(\tau)$,
constructed from the Hankel embedding of logarithmic price returns,
provides a natural dynamical representation of financial markets,
with the observed price series being a one-dimensional projection
of its spectral decomposition. Decomposing $\K(\tau)$ relative to
the Marchenko-Pastur (MP) bulk boundary $\lp$, we show that the integrated
covariance matrix $\X(\tau)=\int_{0}^{\tau}\K(\tau')d\tau'$ exhibits
two distinct diffusive regimes: structural modes diffuse with exponent
$\beta_{\mathrm{s}}\approx1.0$, which indicates near Brownian motion,
and bulk modes that are anomalously superdiffusive with $\beta_{b}\in[1.3,\,1.7]$
are consistent across two intraday timescales. This spectral separation
is explained by a generalized matrix Langevin equation (GLE) with
a quasi-block-diagonal memory kernel, the structural block has a finite
memory integral (finite effective friction), while the bulk block
carries a divergent power-law memory. A third universal exponent $\alpha\approx2$
governs the power-law distribution of inter-regime waiting times.
All three exponents, together with an intrinsic embedding dimension
$\Lopt$ and observation step $\dtopt$, are consistent across two
intraday timescales (M5 and H1) differing by a factor of 12 in physical
time, and across 14 assets spanning equities, indices, foreign exchange,
commodities and fixed-income markets. This supports operator-space
dynamics as a natural level of description for financial markets and
identify the universality class of financial spectral phase transitions. 
\end{abstract}
\noindent \textbf{Keywords:} econophysics; random matrix theory; anomalous
diffusion; generalized Langevin equation; financial markets; spectral
phase transitions; Marchenko-Pastur law.

\section{Introduction}

Financial price series exhibit well-documented stylised facts that
resist a unified description, including persistent intra-regime patterns,
abrupt structural breaks, volatility clustering, non-stationarity,
and anomalous diffusion of returns on multiple timescales~\cite{cont2001,bouchaud2003,mantegna1999}.
Existing frameworks address these features in isolation: GARCH models
capture volatility clustering~\cite{engle1982,bollerslev1986}; Hidden
Markov Models describe regime switching~\cite{hamilton1989}; Random
Matrix Theory (RMT) filters noise from \emph{cross-sectional} covariance
matrices~\cite{laloux1999,plerou1999}. None provides a single dynamical
equation from which all observed phenomena emerge simultaneously.

The root cause of this fragmentation is a modelling choice, not a
fundamental limitation, all these approaches treat the scalar price
(or return) series $P(t)$ as the primary dynamical object. Here we
take a different stance. We identify the \emph{temporal covariance
operator} $\K(\tau)\in\mathbb{R}^{L\times L}$, constructed from the
Hankel embedding~\cite{broomhead1986,takens1981} of a single return
series within a sliding window indexed by slow time $\tau$, as an
effective dynamical object capturing the evolving correlation structure.
The observed return at internal time $t$ within the embedding window
centred at slow time $\tau$ admits the spectral reconstruction 
\begin{equation}
P(t;\tau)=\underbrace{\sum_{k:\,\lambda_{k}(\tau)>\lambda_{+}(\tau)}a_{k}(\tau)\,\phi_{k}(t;\tau)}_{\text{structural signal}}+\underbrace{\sum_{k:\,\lambda_{k}(\tau)\leq\lambda_{+}(\tau)}a_{k}(\tau)\,\phi_{k}(t;\tau)}_{\text{bulk noise}},\label{eq:decomp}
\end{equation}
where $t\in\{1,\ldots,L\}$ is the internal (fast) lag index within
the embedding window---distinct from the slow time $\tau$ that indexes
successive windows---$\{\lambda_{k}(\tau),\phi_{k}(t;\tau)\}$ are
the eigenpairs of $\mathbf{K}(\tau)$, $\lambda_{+}(\tau)$ is the
Marchenko--Pastur upper boundary, and 
\begin{equation}
a_{k}(\tau)=\sum_{t=1}^{L}P(t;\tau)\,\phi_{k}(t;\tau)\label{eq:scores}
\end{equation}
are the projection scores of the return series onto the $k$-th eigenmode
within window $\tau$. This decomposition is directly analogous to
the distinction between a density matrix, the fundamental object of
quantum mechanics, and the scalar expectation values accessible to
any particular measurement.

The apparent complexity of $P(t)$, including long memory, non-stationarity,
and volatility clustering, can be interpreted as arising from projecting
a high-dimensional operator process onto a one-dimensional observable.
Our main result is that $\K(\tau)$ obeys a \emph{generalized matrix
Langevin equation} (GLE) with a quasi-block-diagonal memory kernel,
predicting two distinct diffusion exponents for the structural and
bulk subspaces, together with a universal first-passage exponent for
spectral phase transitions. Empirical analysis across 14 assets and
five asset classes confirms three universal exponents: $(\beta_{\mathrm{s}},\,\beta_{\mathrm{b}},\,\alpha)$.

\section{Construction of the Covariance Operator}

% =======================================================

Let $P(t)=\ln[S(t)/S(t-1)]-\mu(\tau)$ be the mean-centred log-return,
where $\mu(\tau)$ is the sample mean within a window of length $W$
centred at slow time $\tau$. Fix an embedding dimension $L<W$ and
define the Hankel matrix 
\begin{equation}
\Hmat(\tau)\in\mathbb{R}^{N\times L},\quad H_{ki}=P(\tau-W+k+i-1),\quad N=W-L.\label{eq:hankel}
\end{equation}
The empirical covariance operator is 
\begin{equation}
\K(\tau)=\frac{1}{N-1}\bigl[\Hmat(\tau)^{\top}\Hmat(\tau)-N\bar{h}\bar{h}^{\top}\bigr]\in\mathbb{R}^{L\times L},\label{eq:K}
\end{equation}
with $\bar{h}$ the column-mean vector. Eigendecomposition yields
$\K(\tau)\varphi_{k}(\tau)=\lambda_{k}(\tau)\varphi_{k}(\tau)$, ordered
$\lambda_{1}\geq\cdots\geq\lambda_{L}$. The MP upper bulk boundary
is 
\begin{equation}
\lp(\tau)=\hat{\sigma}^{2}(\tau)\bigl(1+\sqrt{q}\bigr)^{2},\quad q=L/N,\label{eq:lp}
\end{equation}
with $\hat{\sigma}^{2}(\tau)$ the empirical variance of $\Hmat(\tau)$.
Modes with $\lambda_{k}>\lp$ carry structural information, the remainder
constitute the bulk. The instantaneous structural dimension is $m(\tau)=\#\{k:\lambda_{k}(\tau)>\lp(\tau)\}$.

The structural and bulk projectors are $\Ps(\tau)=\sum_{k:\lambda_{k}>\lp}\varphi_{k}\varphi_{k}^{\top}$
and $\Pb(\tau)=\mathbf{I}-\Ps(\tau)$. We define the structural and
bulk components of the integrated process 
\begin{equation}
\X(\tau)=\sum_{\tau'<\tau}\K(\tau')\,\Delta\tau\label{eq:X}
\end{equation}
, where $\Delta\tau$ is the time between observations, and $\Xs=\X\Ps$,
$\Xb=\X\Pb$. The mean-square displacement (MSD) of each component,
\begin{equation}
\mathrm{MSD}_{\mathrm{s,b}}(s)=\E\!\left[\Fnorm{\X_{\mathrm{s,b}}(\tau+s)-\X_{\mathrm{s,b}}(\tau)}\right]\sim s{}^{\beta_{\mathrm{s,b}}},\label{eq:msd}
\end{equation}
, where $F$ stands for Frobenius norm, defines the diffusion exponents
in the asymptotic regime.

\section{Generalized Matrix Langevin Equation}

\subsection{Equation of motion}

We model the evolution of $\K(\tau)$ on the slow timescale by the
matrix GLE 
\begin{equation}
\frac{\Delta\K}{\Delta\tau}=-\sum_{\tau'<\tau}\boldsymbol{\Gamma}(\tau-\tau')\,\K(\tau')\,\Delta\tau+\boldsymbol{\xi}(\tau)+\delta\K\,\Delta N_{\tau},\label{eq:GLE}
\end{equation}
where $\boldsymbol{\Gamma}(\tau)$ is the matrix memory kernel, $\boldsymbol{\xi}(\tau)$
is coloured noise related to the kernel by the fluctuation-dissipation
theorem (FDT)~\cite{kubo1966,mori1965,zwanzig2001} 
\begin{equation}
\E\!\left[\boldsymbol{\xi}(\tau)\boldsymbol{\xi}^{\top}(0)\right]=\|\mathbf{K}\|_{F}^{2}\,\boldsymbol{\Gamma}(\tau),\label{eq:fdt}
\end{equation}
where $\|\mathbf{K}\|_{F}^{2}=\mathrm{Tr}(\mathbf{K}^{\top}\mathbf{K})$
is the squared Frobenius norm of $\mathbf{K}$, playing the role of
an effective temperature that sets the fluctuation amplitude proportional
to the current spectral energy of the operator, and $\Delta N_{\tau}$
is a Poisson process for large exogenous shocks. The noise is coloured,
not white Gaussian, it carries the same temporal structure as the
kernel, and the two are linked by the FDT as a single object.

The GLE in Eq. (\ref{eq:GLE}) is not introduced as a phenomenological
ansatz but follows from a Mori--Zwanzig projection onto the operator
subspace, see Supplementary Material (SM) Section S1 for the explicit
construction. As demonstrated in SM, the memory kernel $\Gamma(\tau)$
and the random force $\boldsymbol{\xi}(\tau)$ are coupled via a matrix-valued
Fluctuation-Dissipation Theorem (FDT, Eq(\ref{eq:fdt})), ensuring
consistency with causality and energy balance in the Frobenius metric. 

\subsection{Quasi-block-diagonal memory kernel}

To account for the coexistence of $\beta_{\mathrm{s}}\neq\beta_{\mathrm{b}}$,
the kernel must carry different memory structures in the two spectral
subspaces. In the instantaneous eigenbasis of $\K$ we propose that
\begin{equation}
\boldsymbol{\Gamma}(\tau)\approx\begin{pmatrix}\Gs(\tau) & \boldsymbol{\varepsilon}(\tau)\\
\boldsymbol{\varepsilon}^{\top}(\tau) & \Gb(\tau)
\end{pmatrix},\quad\|\boldsymbol{\varepsilon}\|\ll\|\Gs\|,\|\Gb\|.\label{eq:kernel}
\end{equation}
The weak off-diagonal block $\boldsymbol{\varepsilon}$ couples the
two subspaces and induces a small observed excess $\beta_{\mathrm{s}}-1>0$.

For the \emph{structural block}, the key condition is a convergent
memory integral~\cite{morgado2002}: 
\begin{equation}
\gamma_{s}^{\mathrm{eff}}=\sum_{n=0}^{\infty}\left\Vert \boldsymbol{\Gamma}_{s}(n\,\Delta\tau)\right\Vert _{F}\Delta\tau<\infty.\label{eq:gamma_eff}
\end{equation}
A finite, nonzero Frobenius-norm integral is both necessary and sufficient
for normal diffusion ($\beta_{s}=1$) in the asymptotic regime~\cite{morgado2002},
\emph{regardless} of whether the memory itself is short or long-ranged.
The detailed shape of $\Gs(\tau)$ at intermediate times does not
matter; only its integral does. In the asymptotic limit, $\Gs$ acts
as an effective delta-function friction $\gamma_{\mathrm{s}}^{\mathrm{eff}}\,\delta(\tau)$.

For the \emph{bulk block}, we assume a power-law memory: 
\begin{equation}
\left\Vert \boldsymbol{\Gamma}_{b}(\tau)\right\Vert _{F}\sim\tau^{-(2-\nu)},\quad0<\nu<1,\label{eq:bulk_kernel}
\end{equation}
whose Frobenius-norm integral diverges. A divergent integral with
this power-law form produces anomalous superdiffusion $\beta_{\mathrm{b}}=1+\nu$~\cite{morgado2002}.
Physically, the $L\gg m(\tau)$ bulk modes are quasi-degenerate and
strongly coupled, their collective stochastic mixing via Dyson-type
eigenvalue dynamics~\cite{dyson1962,anderson2010} generates effective
long-range correlations in the bulk sector even though each individual
mode is locally Markovian.

The diffusion exponents $\beta_{s}$ and $\beta_{b}$ characterize
the scaling of the integrated operator $\boldsymbol{X}(\tau)=\int_{0}^{\tau}\boldsymbol{K}(s)ds$.
By projecting the matrix dynamics onto the Frobenius norm $\|\cdot\|_{F}$,
we establish (see SM, Section S2) that $\beta$ is uniquely determined
by the integrability of the memory kernel $\boldsymbol{\Gamma}(\tau)$. 

\subsection{Asymptotic diffusion exponent}

Following~\cite{morgado2002}, $\beta_{\mathrm{s,b}}$ is estimated
by log-log regression of $\mathrm{MSD}_{\mathrm{s,b}}(s)$ restricted
to the upper 70\% of available lags. Discarding the lower 30\% where
transient behaviour dominates.

The small excess $\beta_{\mathrm{s}}>1$ (Table~\ref{tab:exponents})
is attributed to three sources in decreasing order of expected magnitude:
(i)~weak inter-block coupling $\boldsymbol{\varepsilon}$, which
injects a superdiffusive component from the bulk; (ii)~the Poisson
term $\delta\K\,\Delta N_{\tau}$, contributing sporadic large increments;
and (iii)~finite-sample regression uncertainty.

\subsection{Fluctuation--dissipation balance}

Applying It�'s lemma to the Frobenius energy $E=\|\mathbf{K}\|_{F}^{2}=\tr(\mathbf{K}^{\top}\mathbf{K})$,
the increment decomposes as 
\begin{equation}
\Delta E=\underbrace{2\,\tr(\mathbf{K}\,\Delta\mathbf{K}^{\top})}_{\text{linear term}}+\underbrace{\tr(\Delta\mathbf{K}\,\Delta\mathbf{K}^{\top})}_{\text{quadratic term}},\label{eq:ito_energy}
\end{equation}
with no approximation. The linear term is odd in $\Delta\mathbf{K}$
and changes sign with the direction of motion of $\mathbf{K}$; near
stationarity it acts as an energy sink. The quadratic term is positive
definite and acts as a source. We therefore define 
\begin{align}
D(\tau) & =-2\,\tr\!\bigl(\mathbf{K}(\tau)\,\Delta\mathbf{K}(\tau)^{\top}\bigr),\label{eq:dissip}\\
I(\tau) & =\phantom{-}\tr\!\bigl(\Delta\mathbf{K}(\tau)\,\Delta\mathbf{K}(\tau)^{\top}\bigr),\label{eq:inject}
\end{align}
as the instantaneous dissipation and injection rates, and the energy-balance
ratio as 
\begin{equation}
R_{\mathrm{FDT}}=\frac{\langle D\rangle}{\langle I\rangle},\label{eq:R_fdt}
\end{equation}
where $\langle\cdot\rangle$ denotes a time average. At stationarity,
$\langle\Delta E\rangle=0$ implies $\langle D\rangle=\langle I\rangle$
and hence $R_{\mathrm{FDT}}=1$ exactly, by Eq.~(\ref{eq:ito_energy}).
This result follows solely from the It� decomposition and holds independently
of the diffusive regime: in particular, it does not require the effective
friction to be finite, and remains valid in the superdiffusive regime
where the effective friction vanishes. Deviations of $R_{\mathrm{FDT}}$
from unity therefore signal departures from stationarity rather than
transitions between dynamical regimes. Empirically, all liquid assets
yield $R_{\mathrm{FDT}}\approx1.0$ to within the measurement uncertainty
(Table~\ref{tab:exponents}), confirming that the operator $K(\tau)$
operates near a stationary state over the observation window, consistently
with the stationarity assumption underlying the Hankel embedding.

\section{Spectral Phase Transitions}

\subsection{Regime changes as first-passage events}

Define the spectral gap $G(\tau)=\lambda_{m(\tau)}(\tau)-\lambda_{+}(\tau)$:
the signed distance between the marginal structural eigenvalue and
the Marchenko--Pastur boundary. A regime change (increment or decrement
of $m(\tau)$) occurs when $G(\tau)=0$, so the sequence of inter-regime
durations $\{\tau_{k}^{*}\}$ is precisely the sequence of first-passage
times of $G$ through zero. Applying the GLE eigenvalue calculus~\cite{dyson1962,anderson2010},
the gap evolves as 
\begin{equation}
\Delta G=\mu(G,\tau)\,\Delta\tau_{\mathrm{opt}}+\sigma(G,\tau)\,\Delta W_{\tau},
\end{equation}
with $\mu(G)=a(\tau)G+\mathcal{O}(G^{2})$ and $\sigma(0,\tau)>0$.
As $G\to0$, the drift is linear while the diffusion coefficient remains
finite, which places the transition in the universality class of one-dimensional
diffusions at a regular critical point~\cite{redner2001,majumdar1999}.

For such processes, first-passage theory predicts a power-law distribution
of return times $g(\tau^{*})\sim(\tau^{*})^{-\alpha}$ with the exponent
controlled by the local drift near the boundary \cite{redner2001,majumdar1999}.
In the marginal regime where the drift vanishes, values near $\alpha\approx2$
are expected. We therefore treat $\alpha$ as an empirical exponent,
and verify that the measured values$\hat{\alpha}\in[1.87,\,2.45]$
are consistent with this scenario.

%whose exponent  depends on the sign and magnitude of the drift coefficient . In the symmetric case  (driftless diffusion), ; for  (drift toward the boundary), an exponential cutoff develops; for  (drift away from the boundary),  is recovered in the marginal regime. Real markets are plausibly near this marginal regime on average, since structural eigenvalues spend comparable time above and below  across different market conditions. This heuristic argument suggests  as a natural value, consistent with interface depinning and record-breaking statistics, but a rigorous derivation from the GLE memory kernel remains an open problem. We therefore treat  as an empirically measured exponent, calibrated against the AR(1) family as described in Section, and note that the measured values  are consistent with this picture without requiring it.

\subsection{\label{subsec:Adiabatic-eigenvector-dynamics}Adiabatic eigenvector
dynamics}

When $\mathbf{K}(\tau)$ varies slowly in $\tau$, differentiating
the eigenvalue equation $\mathbf{K}(\tau)\boldsymbol{\phi}_{k}(\cdot;\tau)=\lambda_{k}(\tau)\boldsymbol{\phi}_{k}(\cdot;\tau)$
with respect to $\tau$ and projecting onto $\boldsymbol{\phi}_{j}$
($j\neq k$) gives 
\begin{equation}
\partial_{\tau}\boldsymbol{\phi}_{k}=\sum_{j\neq k}\frac{\boldsymbol{\phi}_{j}^{\top}(\partial_{\tau}\mathbf{K})\boldsymbol{\phi}_{k}}{\lambda_{k}-\lambda_{j}}\,\boldsymbol{\phi}_{j},\label{eq:adiabatic}
\end{equation}
the discrete analogue of adiabatic perturbation theory for self-adjoint
operators~\cite{kato1995,messiah1961}. Here the diagonal term $j=k$
vanishes by differentiation of the normalisation condition $\|\boldsymbol{\phi}_{k}\|^{2}=1$.
Equation~\eqref{eq:adiabatic} applies to structural modes ($\lambda_{k}>\lambda_{+}$);
eigenvectors within the MP bulk are stochastically mixed and do not
admit a smooth adiabatic trajectory. As $\lambda_{k}\to\lambda_{j}$,
the denominator vanishes and the eigenvectors rotate rapidly---the
spectral precursor of a regime change. The order parameter of the
transition is 
\begin{equation}
\Delta(\tau)=\min_{k\neq j}|\lambda_{k}(\tau)-\lambda_{j}(\tau)|.\label{eq:order_param}
\end{equation}


\section{Bulk Spectral Universality}

In the quasi-degenerate bulk subspace, stochastic mixing via~(\ref{eq:adiabatic})
induces Brownian rotation of eigenvectors. Projections of returns
onto bulk modes become approximately Gaussian, and the spectral density
converges to the MP law~\cite{marchenko1967,bai2010}: 
\begin{equation}
\rho_{\mathrm{MP}}(\lambda)=\frac{\sqrt{(\lambda_{+}-\lambda)(\lambda-\lambda_{-})}}{2\pi q\hat{\sigma}^{2}\,\lambda},\label{eq:mp}
\end{equation}
where $\lambda_{\pm}=\hat{\sigma}^{2}(1\pm\sqrt{q})^{2}$.

The observed relative $L^{2}$ error of $\approx3.7\%$ between the
empirical bulk histogram and~(\ref{eq:mp}) is attributed to short-range
spurious correlations introduced by the overlapping Hankel window
construction, which slightly deforms the bulk spectrum. This systematic
offset is a finite-size construction artefact and does not affect
the structural eigenvalue dynamics.

\section{Empirical Results}

We analyse log-returns for 14 assets across equities, fixed-income,
foreign exchange, commodity and rates markets (Table~\ref{tab:exponents})
at two intraday timescales: M5 (5-minute bars) and H1 (1-hour bars),
each spanning 5~years. We choose the optimal parameters according
to a robustness criterion defined in subsection (\ref{subsec:Intrinsic-embedding-dimension}).

\subsection{Bulk spectral density (Fig.~\ref{fig:mp})}

\begin{figure}[htbp]
\centering \includegraphics{PETR4_M5_fig1_mp} \caption{\textbf{Bulk eigenvalue density, PETR4 (M5).} Aggregated, normalised
bulk eigenvalue histogram (blue bars). Red curve: theoretical MP density
Eq.~(\ref{eq:mp}) with $q=0.25$, $\hat{\sigma}^{2}=1$. Vertical
dashed lines: $\lambda_{\pm}$. \emph{Inset}: relative $L^{2}$ error
between empirical bulk histogram and MP law. The systematic $\approx3.7\%$
deviation is attributed to short-range spurious correlations from
the overlapping Hankel window construction.}
\label{fig:mp}
\end{figure}

Figure~\ref{fig:mp} shows the aggregated, normalised bulk eigenvalue
histogram for PETR4 (M5 timescale). The relative $L^{2}$ error to
the theoretical MP density~(\ref{eq:mp}) is $3.7\%$ (inset). The
small systematic deviation is attributed to Hankel window overlap.

\subsection{Structural diffusion exponent $\beta_{\mathrm{s}}$ (Fig.~\ref{fig:exponents}a,c)}

\begin{figure}[htbp]
\centering \includegraphics{PETR4_fig2_exponents} \caption{\textbf{Universal exponents, PETR4.} (a)~Log-log MSD of $\Xs(\tau)$
vs.\ lag $s$ (in units of $\dtopt$) at M5 (blue) and H1 (red).
Dashed: reference slope $\beta=1$. Asymptotic exponent $\bstruct\approx1.0$
estimated over the upper 70\% of lags~\cite{morgado2002}. (b)~Log-log
histogram of inter-regime durations $\tau^{*}$ at M5; dashed: slope
$\alpha=2$. (c)~$\bstruct$ vs.\ observation step $\Delta\tau_{0}$;
green band: optimal range $\dtopt\approx15$-$25$~candles. (d)~$\hat{\alpha}$
vs.\ $\Delta\tau_{0}$; dashed: $\alpha=2$. }
\label{fig:exponents}
\end{figure}

Figure~\ref{fig:exponents}(a) shows the log-log MSD of $\Xs(\tau)$
for PETR4 at M5 and H1. Both timescales yield $\bstruct\approx1.0$,
nearly a Brownian motion. The exponent is estimated over the upper
70\% of lags~\cite{morgado2002}. The small excess $\delta\beta_{\mathrm{s}}=\bstruct-1>0$
is consistent with weak bulk-structure coupling and the Poisson term
(Sec.~III.C).

Figure~\ref{fig:exponents}(c) shows $\bstruct$ vs.\ observation
step $\Delta\tau_{0}$. For $\Delta\tau_{0}<\dtopt$, overlapping
successive windows bias the MSD upward; for $\Delta\tau_{0}>\dtopt$,
M5 and H1 decouple into their respective physical asymptotic regimes,
breaking inter-scale coherence. The optimal range $\dtopt\approx15$-$25$~candles
is the regime where successive increments are approximately uncorrelated
and both timescales share the same universal dynamics.

\subsection{Bulk diffusion exponent $\hat{\beta}_{b}$}

The bulk MSD exponent $\hat{\beta}_{b}$ is superdiffusive for all
assets and both timescales (Table~\ref{tab:exponents}), with values
in the range $[1.3,\,1.7]$. This spread reflects genuine differences
in the collective eigenvalue dynamics of the bulk sector across asset
classes and sampling frequencies: assets with denser microstructure
or higher liquidity tend toward higher values of $\hat{\beta}_{b}$,
consistent with stronger Dyson-type stochastic mixing among quasi-degenerate
bulk modes (Sec.~\ref{subsec:Adiabatic-eigenvector-dynamics}). The
central tendency $\hat{\beta}_{b}\approx1.5$, is representative of
the liquid equity and FX futures sample; the AR(1) reference processes
yield $\hat{\beta}_{b}\approx\hat{\beta}_{s}\approx1$ for all $\varphi$
(Supplementary Table~1), confirming that the structural--bulk exponent
separation is a genuine dynamical signature of real markets and not
an artefact of the Hankel embedding.

\subsection{Inter-regime waiting times $\alpha$ (Fig.~\ref{fig:exponents}b,d)}

Figure~\ref{fig:exponents}(b) shows the log-log histogram of inter-regime
intervals $\tau^{*}$ for PETR4 at M5, where a power-law tail with
$\hat{\alpha}\approx2$ is observed. Figure~\ref{fig:exponents}(d)
shows $\hat{\alpha}$ vs.\ $\Delta\tau_{0}$; convergence to $\approx2$
is robust for both timescales at smaller $\Delta\tau_{0}$. Exponential
and Gaussian waiting-time distributions, characteristic of Poisson
and GARCH models, are strongly rejected.

\subsection{\label{subsec:Intrinsic-embedding-dimension}Intrinsic embedding
dimension $\Lopt$ (Fig.~\ref{fig:lopt})}

\begin{figure}[htbp]
\centering \includegraphics{PETR4_fig3_Lopt} \caption{\textbf{Intrinsic embedding dimension, PETR4.} $\bstruct$ vs.\ $L$
for $n=W/L\in\{4,5,6\}$ at M5 (cool colours) and H1 (warm colours).
Yellow band: $\Lopt\in[70,90]$, defined by the inter-scale coherence
criterion~(\ref{eq:L_opt}). Both timescales converge in the shaded
region, identifying $\Lopt\approx80$ as an intrinsic scale invariant
under a factor-of-12 change in physical time.}
\label{fig:lopt}
\end{figure}

Figure~\ref{fig:lopt} shows $\bstruct$ vs.\ $L$ for $n=W/L\in\{4,5,6\}$
and both timescales. For $L<\Lopt$, the spectral projector $\Ps$
is geometrically unstable, the spectral gap is insufficient to isolate
structural eigenmodes, producing a subdiffusive artefact in H1 timescale
$\bstruct<1$. For $L>\Lopt$, lags beyond the process correlation
time $\tau_{c}$ contaminate $\Ps$ with bulk eigenmodes. 

The parameters $(\Delta\tau_{\mathrm{opt}},\,L_{\mathrm{opt}})$ are
selected by a cross-scale robustness criterion that requires no theoretical
anchor. Both optimizations share the same composite score structure:
a local-stability term and an inter-scale coherence term, each normalized
by its grid average to render them dimensionless and commensurable
without any free parameter.

For $\Delta\tau_{\mathrm{opt}}$, we compute $\hat{\beta}_{\mathrm{struct}}$
at both M5 and H1 for each candidate step and define the composite
score at each interior grid point (excluding the two outermost points
on each side) as 
\begin{equation}
\mathcal{S}(\Delta\tau_{0})=\frac{\bar{\sigma}_{\mathrm{local}}^{2}(\Delta\tau_{0})}{\overline{\bar{\sigma}_{\mathrm{local}}^{2}}}+\frac{B(\Delta\tau_{0})}{\bar{B}},\label{eq:score_dtau}
\end{equation}
where 
\begin{equation}
\bar{\sigma}_{\mathrm{local}}^{2}(\Delta\tau_{0})=\frac{1}{2}\sum_{f\,\in\,\{\mathrm{M5,H1}\}}\mathrm{Var}\!\left[\hat{\beta}_{\mathrm{struct}}^{(f)}(\Delta\tau_{0}+jh),\;j\in\{-2,-1,0,1,2\}\right]\label{eq:termA_dtau}
\end{equation}
is the mean of the local variances computed independently for each
timescale over the five-point window centered at $\Delta\tau_{0}$
(grid step $h$), and 
\begin{equation}
B(\Delta\tau_{0})=\left|\hat{\beta}_{\mathrm{struct}}^{(\mathrm{M5})}(\Delta\tau_{0})-\hat{\beta}_{\mathrm{struct}}^{(\mathrm{H1})}(\Delta\tau_{0})\right|\label{eq:termB_dtau}
\end{equation}
is the inter-scale discrepancy. The grid averages $\overline{\bar{\sigma}_{\mathrm{local}}^{2}}$
and $\bar{B}$ are computed over all interior points. The optimal
lag is 
\begin{equation}
\Delta\tau_{\mathrm{opt}}=\operatorname*{arg\,min}_{\Delta\tau_{0}}\;\mathcal{S}(\Delta\tau_{0}).\label{eq:dtau_opt}
\end{equation}
The first term selects the plateau of each curve independently, penalizing
regions where $\hat{\beta}_{\mathrm{struct}}$ is still varying, while
the second penalizes systematic inter-scale discrepancy. Separating
the two contributions prevents pathological cases where the curves
cross at small $\Delta\tau_{0}$ (driving $B\to0$ artificially) or
where they are parallel and offset (making $B$ uninformative). The
five-point window (second neighbors) provides robustness against isolated
noise spikes that would corrupt a three-point estimate.

For $L_{\mathrm{opt}}$, an analogous composite score is defined,
with the sole difference that only first neighbors are used in the
local-stability term, reflecting the coarser resolution of the embedding-dimension
grid: 
\begin{equation}
\mathcal{S}_{L}(L)=\frac{\bar{\sigma}_{L,\mathrm{local}}^{2}(L)}{\overline{\bar{\sigma}_{L,\mathrm{local}}^{2}}}+\frac{B_{L}(L)}{\bar{B}_{L}},\label{eq:score_L}
\end{equation}
where 
\begin{equation}
\bar{\sigma}_{L,\mathrm{local}}^{2}(L)=\frac{1}{2}\sum_{f\,\in\,\{\mathrm{M5,H1}\}}\mathrm{Var}\!\left[\hat{\beta}_{\mathrm{struct}}^{(f)}(L+jh_{L}),\;j\in\{-1,0,1\}\right]\label{eq:termA_L}
\end{equation}
is the mean of the local variances over the three-point window (grid
step $h_{L}$), and 
\begin{equation}
B_{L}(L)=\left|\hat{\beta}_{\mathrm{struct}}^{(\mathrm{M5})}(L)-\hat{\beta}_{\mathrm{struct}}^{(\mathrm{H1})}(L)\right|\label{eq:termB_L}
\end{equation}
is the inter-scale discrepancy at fixed multiplier. The optimal embedding
dimension is 
\begin{equation}
L_{\mathrm{opt}}=\operatorname*{arg\,min}_{L}\;\mathcal{S}_{L}(L).\label{eq:L_opt}
\end{equation}
At $L_{\mathrm{opt}}$, the score is by construction at its minimum
with respect to both local stability and inter-scale coherence, the
structural exponent is by construction insensitive to the choice of
window multiplier within the range $\{4L,\,5L,\,6L\}$ at that point,
and we adopt the window of $5L$ candles as the reference throughout.
The two optimizations are performed in series: $\Delta\tau_{\mathrm{opt}}$
is determined first with $L=70$ fixed, and $L_{\mathrm{opt}}$ is
then determined with $\Delta\tau_{\mathrm{opt}}$ fixed.

\subsection{Universal exponents across asset classes (Table~\ref{tab:exponents})}

\begin{table}
\centering \caption{\textbf{Universal exponents across asset classes and markets.} $\bstruct$:
structural MSD exponent; $\hat{\alpha}$: inter-regime waiting-time
exponent; $\bbulk$: bulk MSD exponent; $\Lopt$: optimal embedding
dimension (timescale-invariant, single column). M5: 5-min bars; H1:
1-hour bars. Asset classes: \fcolorbox{black}{yellow!30}{\rule{0pt}{0.25em}\rule{0.25em}{0pt}}
Equities, \fcolorbox{black}{blue!20}{\rule{0pt}{0.25em}\rule{0.25em}{0pt}}
Index, \fcolorbox{black}{orange!30}{\rule{0pt}{0.25em}\rule{0.25em}{0pt}}
Forex, \fcolorbox{black}{green!20}{\rule{0pt}{0.25em}\rule{0.25em}{0pt}}
Commodities and \fcolorbox{black}{red!20}{\rule{0pt}{0.25em}\rule{0.25em}{0pt}}
Rates.}
\label{tab:exponents}\setlength{\tabcolsep}{2pt}{\small{} }%
\begin{tabular}{|l||ccc||ccc||c|}
\cline{2-7} \cline{3-7} \cline{4-7} \cline{5-7} \cline{6-7} \cline{7-7} 
\multicolumn{1}{l|}{} & \multicolumn{3}{c|}{\textbf{\footnotesize{}M5 (5-min)}} & \multicolumn{3}{c|}{\textbf{\footnotesize{}H1 (1-hour)}} & \multicolumn{1}{c}{}\tabularnewline
\hline 
\hline 
\textbf{\footnotesize{}Asset/Market} & \textbf{\footnotesize{}$\hat{\alpha}$ } & \textbf{\footnotesize{}$\bstruct$ } & \textbf{\footnotesize{}$\bbulk$ } & \textbf{\footnotesize{}$\hat{\alpha}$ } & \textbf{\footnotesize{}$\bstruct$ } & \textbf{\footnotesize{}$\bbulk$ } & \textbf{\footnotesize{}$\Lopt$}\tabularnewline
\hline 
{\footnotesize{}\cellcolor{yellow!30}PETR4/B3 } & {\footnotesize{}$2.023\pm0.181$ } & {\footnotesize{}$1.052\pm0.006$ } & {\footnotesize{}$1.268\pm0.005$ } & {\footnotesize{}$1.849\pm0.122$ } & {\footnotesize{}$1.111\pm0.046$ } & {\footnotesize{}$1.556\pm0.023$ } & {\footnotesize{}80}\tabularnewline
\cline{2-8} \cline{3-8} \cline{4-8} \cline{5-8} \cline{6-8} \cline{7-8} \cline{8-8} 
{\footnotesize{}\cellcolor{yellow!30}ITUB4/B3 } & {\footnotesize{}$2.013\pm0.170$ } & {\footnotesize{}$1.020\pm0.012$ } & {\footnotesize{}$1.603\pm0.008$ } & {\footnotesize{}$1.973\pm0.153$ } & {\footnotesize{}$0.925\pm0.019$ } & {\footnotesize{}$1.546\pm0.018$ } & {\footnotesize{}100}\tabularnewline
\cline{2-8} \cline{3-8} \cline{4-8} \cline{5-8} \cline{6-8} \cline{7-8} \cline{8-8} 
{\footnotesize{}\cellcolor{yellow!30}VALE3/B3 } & {\footnotesize{}$2.247\pm0.247$ } & {\footnotesize{}$0.965\pm0.010$ } & {\footnotesize{}$1.653\pm0.014$ } & {\footnotesize{}$2.245\pm0.281$ } & {\footnotesize{}$0.857\pm0.063$ } & {\footnotesize{}$1.503\pm0.044$ } & {\footnotesize{}120}\tabularnewline
\cline{2-8} \cline{3-8} \cline{4-8} \cline{5-8} \cline{6-8} \cline{7-8} \cline{8-8} 
{\footnotesize{}\cellcolor{yellow!30}AAPL/NYSE } & {\footnotesize{}TBD } & {\footnotesize{}TBD } & {\footnotesize{}TBD } & {\footnotesize{}TBD } & {\footnotesize{}TBD } & {\footnotesize{}TBD } & {\footnotesize{}TBD}\tabularnewline
\cline{2-8} \cline{3-8} \cline{4-8} \cline{5-8} \cline{6-8} \cline{7-8} \cline{8-8} 
{\footnotesize{}\cellcolor{blue!20}S\&P~500/NYSE } & {\footnotesize{}TBD } & {\footnotesize{}TBD } & {\footnotesize{}TBD } & {\footnotesize{}TBD } & {\footnotesize{}TBD } & {\footnotesize{}TBD } & {\footnotesize{}TBD}\tabularnewline
\cline{2-8} \cline{3-8} \cline{4-8} \cline{5-8} \cline{6-8} \cline{7-8} \cline{8-8} 
{\footnotesize{}\cellcolor{blue!20}Nikkei~225/TSE} & {\footnotesize{}TBD } & {\footnotesize{}TBD } & {\footnotesize{}TBD } & {\footnotesize{}TBD } & {\footnotesize{}TBD } & {\footnotesize{}TBD } & {\footnotesize{}TBD}\tabularnewline
\cline{2-8} \cline{3-8} \cline{4-8} \cline{5-8} \cline{6-8} \cline{7-8} \cline{8-8} 
{\footnotesize{}\cellcolor{blue!20}CAC~40/XPAR} & {\footnotesize{}TBD } & {\footnotesize{}TBD } & {\footnotesize{}TBD } & {\footnotesize{}TBD } & {\footnotesize{}TBD } & {\footnotesize{}TBD } & {\footnotesize{}TBD}\tabularnewline
\cline{2-8} \cline{3-8} \cline{4-8} \cline{5-8} \cline{6-8} \cline{7-8} \cline{8-8} 
{\footnotesize{}\cellcolor{blue!20}WIN\$N/B3} & {\footnotesize{}$2.254\pm0.262$ } & {\footnotesize{}$1.015\pm0.007$ } & {\footnotesize{}$1.640\pm0.012$ } & {\footnotesize{}$2.312\pm0.268$ } & {\footnotesize{}$1.019\pm0.041$ } & {\footnotesize{}$1.576\pm0.030$ } & {\footnotesize{}70}\tabularnewline
\cline{2-8} \cline{3-8} \cline{4-8} \cline{5-8} \cline{6-8} \cline{7-8} \cline{8-8} 
{\footnotesize{}\cellcolor{orange!30}EUR/USD/OTC} & {\footnotesize{}TBD } & {\footnotesize{}TBD } & {\footnotesize{}TBD } & {\footnotesize{}TBD } & {\footnotesize{}TBD } & {\footnotesize{}TBD } & {\footnotesize{}TBD}\tabularnewline
\cline{2-8} \cline{3-8} \cline{4-8} \cline{5-8} \cline{6-8} \cline{7-8} \cline{8-8} 
{\footnotesize{}\cellcolor{orange!30}WDO\$N/B3} & {\footnotesize{}$2.240\pm0.225$ } & {\footnotesize{}$1.031\pm0.005$ } & {\footnotesize{}$1.681\pm0.014$ } & {\footnotesize{}$2.184\pm0.234$ } & {\footnotesize{}$0.719\pm0.066$ } & {\footnotesize{}$1.377\pm0.060$ } & {\footnotesize{}80}\tabularnewline
\cline{2-8} \cline{3-8} \cline{4-8} \cline{5-8} \cline{6-8} \cline{7-8} \cline{8-8} 
{\footnotesize{}\cellcolor{green!20}Gold/OTC} & {\footnotesize{}TBD } & {\footnotesize{}TBD } & {\footnotesize{}TBD } & {\footnotesize{}TBD } & {\footnotesize{}TBD } & {\footnotesize{}TBD } & {\footnotesize{}TBD}\tabularnewline
\cline{2-8} \cline{3-8} \cline{4-8} \cline{5-8} \cline{6-8} \cline{7-8} \cline{8-8} 
{\footnotesize{}\cellcolor{green!20}WTI/CME} & {\footnotesize{}TBD } & {\footnotesize{}TBD } & {\footnotesize{}TBD } & {\footnotesize{}TBD } & {\footnotesize{}TBD } & {\footnotesize{}TBD } & {\footnotesize{}TBD}\tabularnewline
\cline{2-8} \cline{3-8} \cline{4-8} \cline{5-8} \cline{6-8} \cline{7-8} \cline{8-8} 
{\footnotesize{}\cellcolor{red!20}DI1/B3} & {\footnotesize{}$2.116\pm0.197$ } & {\footnotesize{}$1.051\pm0.011$ } & {\footnotesize{}$1.370\pm0.021$ } & {\footnotesize{}$1.982\pm0.160$ } & {\footnotesize{}$1.059\pm0.030$ } & {\footnotesize{}$1.333\pm0.050$ } & {\footnotesize{}30}\tabularnewline
\cline{2-8} \cline{3-8} \cline{4-8} \cline{5-8} \cline{6-8} \cline{7-8} \cline{8-8} 
{\footnotesize{}\cellcolor{red!20}US 10Y/CME} & {\footnotesize{}TBD } & {\footnotesize{}TBD } & {\footnotesize{}TBD } & {\footnotesize{}TBD } & {\footnotesize{}TBD } & {\footnotesize{}TBD } & {\footnotesize{}TBD}\tabularnewline
\hline 
\multicolumn{1}{|l|}{{\footnotesize{}Mean }} & {\footnotesize{}TBD } & {\footnotesize{}TBD } & \multicolumn{1}{c|}{{\footnotesize{}TBD }} & {\footnotesize{}TBD } & {\footnotesize{}TBD } & {\footnotesize{}TBD } & {\footnotesize{}TBD}\tabularnewline
\cline{2-8} \cline{3-8} \cline{4-8} \cline{5-8} \cline{6-8} \cline{7-8} \cline{8-8} 
\multicolumn{1}{|l|}{{\footnotesize{}Std }} & {\footnotesize{}TBD } & {\footnotesize{}TBD } & \multicolumn{1}{c|}{{\footnotesize{}TBD }} & {\footnotesize{}TBD } & {\footnotesize{}TBD } & {\footnotesize{}TBD } & {\footnotesize{}TBD}\tabularnewline
\hline 
\end{tabular}
\end{table}

Table~\ref{tab:exponents} summarises $\bstruct$, $\hat{\alpha}$,
$\bbulk$ and $\Lopt$ for all 14 assets at both timescales. The universality
of all three exponents across asset classes, market microstructures
and sampling frequencies constitutes strong evidence for a universal
dynamical class in operator space.

A notable exception is \textsc{DI1} (Brazilian overnight interest-rate
futures), whose optimal embedding dimension $L_{\mathrm{opt}}=30$
is substantially smaller than for equities and FX futures ($L_{\mathrm{opt}}\in[70,120]$,
Table~\ref{tab:exponents}). A smaller $L_{\mathrm{opt}}$ indicates
that the dynamical complexity of $\mathbf{K}(\tau)$ is fully resolved
by a lower-dimensional state space: structural modes are fewer but
more dominant relative to the bulk, which has fewer effective degrees
of freedom. This is consistent with the simpler autocorrelation structure
of interest-rate futures, whose price dynamics are tightly anchored
to monetary policy expectations and exhibit lower intraday complexity
than equity or currency markets.

\section{Discussion}

\subsection{The fundamental object and its projection onto prices}

Classical approaches to financial dynamics, whether rooted in stochastic
calculus~\cite{blackscholes1973,heston1993}, econometric volatility
models~\cite{engle1982,bollerslev1986}, or the statistical physics
of price fluctuations~\cite{mantegna1999,cont2001}, take the price
process $P(\tau)$, or equivalently the return series $r(\tau)$,
as the primary observable. Structure is then inferred by measuring
properties of this scalar time series: its distribution, autocorrelation
function, variance clustering, and spectral density. The present work
proposes a different starting point: the \emph{temporal covariance
operator} $K(\tau)$ is the fundamental dynamical object, and the
observed price return is its spectral projection onto a one-dimensional
observable via Eq.~(\ref{eq:projection}).

This identification has immediate interpretive consequences. GARCH-type
volatility clustering is the scalar footprint of superdiffusive bulk
operator dynamics; regime changes are first-passage events of $G(\tau)$
through zero; and the apparent long memory in scalar returns~\cite{beran1994},
which has long resisted unambiguous interpretation via Hurst exponents
or DFA, is the projection of the superdiffusive bulk onto a one-dimensional
observable, not an intrinsic property of prices. Standard scalar tests
applied directly to $P(t)$ cannot distinguish these operator-level
phenomena. This establishes a clear hierarchy: $K(\tau)$ acts as
a generator, while $P(t)$ corresponds to a projected observable.

This shift is not merely a change of representation: it promotes a
derived statistical quantity to the role of a dynamical field, in
the same spirit that statistical mechanics replaces the trajectory
of individual particles with a density field or an order parameter.
The operator $K(\tau)$ encodes simultaneously the instantaneous correlation
structure of the market, its spectral decomposition into signal and
noise, and its temporal evolution , none of which are directly accessible
from the scalar price series alone.

The specific form of the generalised matrix Langevin equation derived
from this perspective , with a quasi-block-diagonal memory kernel
coupling the structural and bulk subspaces of $K(\tau)$ , is a new
result that emerges directly from the empirical separation of two
distinct diffusive regimes.

\subsection{Why bulk diffusion is more anomalous than the signal}

The result $\bbulk\approx1.3-1.7>\bstruct\approx1.0$ may appear counterintuitive.
The resolution lies in the collective nature of bulk diffusion. The
$L\gg m$ quasi-degenerate bulk modes are strongly coupled through
Eq.~(\ref{eq:adiabatic}): near-degeneracy amplifies their mutual
rotations. This Dyson-type stochastic mixing~\cite{dyson1962} generates
effective long-range memory in the bulk sector , a divergent $\hat{\Gamma}_{\mathrm{b}}(0)$
, even though each individual mode is locally Markovian. By contrast,
the few $m$ structural modes are spectrally isolated, weakly coupled,
and governed by a memory kernel with a finite integral, yielding near-Brownian
diffusion.

\subsection{Connections to established frameworks}

\paragraph{Random matrix theory in finance.}

The application of RMT to financial correlation matrices, pioneered
by Laloux et al.~\cite{laloux1999} and Plerou et al.~\cite{plerou1999},
established that the bulk of the empirical spectrum is well described
by the Marchenko-Pastur law~\cite{marchenko1967}, with a small number
of eigenvalues carrying genuine market information. This programme,
developed extensively in the context of portfolio optimisation and
eigenvalue cleaning~\cite{bouchaud2003,bun2017,potters2020}, treats
the spectral decomposition as a static property of the market. The
present framework extends this programme in two directions. First,
the MP separation is applied locally at each $\tau$, transforming
a single static snapshot into a dynamical observable, whose stability
is characterised by $\alpha$, a new dynamical quantity with no counterpart
in the static RMT literature. Second, the quasi-block-diagonal memory
kernel is the equation of motion that the structural/bulk separation
implies, an object the static RMT programme, by construction, could
not have produced. These are extensions, not contradictions: the static
MP decomposition is recovered as the long-time limit of the local
decomposition, and results on eigenvalue cleaning~\cite{bun2017}
remain fully valid within our framework.

\paragraph{Anomalous diffusion.}

Building on the anomalous diffusion framework developed in Refs.~\cite{morgado2002,oliveira2019},
by analogy with the scalar GLE, where the diffusion exponent $\beta$
is carried by the \emph{position} $x(t)=\int_{0}^{t}v(t')\,\mathrm{d}t'$
rather than by the velocity $v(t)$ itself, the exponent $\beta_{s,b}$
here belongs to the integrated operator process $\mathbf{X}(\tau)=\int_{0}^{\tau}\mathbf{K}(\tau')\,\mathrm{d}\tau'$
rather than to $\mathbf{K}(\tau)$ directly. This analogy is formal---$\mathbf{K}(\tau)$
is a positive-semidefinite covariance operator, not a signed velocity
field---but it is operationally precise: it identifies $\mathbf{X}(\tau)$
as the correct process from which $\beta_{s,b}$ must be estimated
via the log-log slope of $\mathrm{MSD}_{s,b}(s)$, and rules out estimating
$\beta$ directly from the increments of $\mathbf{K}$. This identification,
grounded in the fluctuation-dissipation relations~\cite{kubo1966},
is what makes a direct comparison between $\bstruct$ and $\bbulk$
physically meaningful. A central result of Ref.~\cite{morgado2002}
is that unbiased estimation of $\beta$ requires computing the MSD
of the position, not of the velocity; the estimation procedure adopted
here implements this principle directly, extracting $\beta$ from
the log-log slope of the MSD of the integrated operator over the tail
of the lag distribution. The coexistence of normal diffusion in the
structural subspace and superdiffusion in the bulk, within the same
process and governed by a single memory kernel, is natural in the
Langevin framework~\cite{metzler2000,bouchaud1990} but has not previously
been identified empirically in financial data.

\subsection{Universality class and intrinsic market scales}

The exponent triple $(\bstruct,\,\bbulk,\,\alpha)$ defines the dynamical
signature of the operator-space framework. Across all liquid assets
studied, $\bstruct$ is consistent with normal diffusion ($\bstruct\approx1$),
$\bbulk$ is superdiffusive ($\bbulk\approx1.3-1.7$), and $\alpha$
is concentrated in the range $[1.87,\,2.45]$, close to the value
of 2 associated with interface depinning and record-breaking statistics~\cite{majumdar1999,redner2001,bray2013}.
The optimal parameters $(\Delta\tau_{\mathrm{opt}},\,L_{\mathrm{opt}})$
vary across assets (Table~\ref{tab:params}), reflecting intrinsic
dynamical scales of each market; their variation occurs within a well-defined
range, providing evidence for a universal dynamical class in operator-space
financial econophysics rather than a fixed-point universality in the
strict sense.

To calibrate the position of each asset within this class, we compare
the measured $\alpha$ against the AR(1) reference curve described
in Supplementary Section~\ref{sec:ar1}, which interpolates monotonically
between the white-noise ceiling ($\alpha_{0}\approx2.58$, $\phi=0$)
and the GBM floor ($\alpha_{\mathrm{GBM}}\approx1.87$, $\phi=1$).
All liquid assets fall within this band, spanning $\hat{\varphi}_{\mathrm{eff}}\in[0.1,\,0.9]$.
Most cluster in the range $\hat{\varphi}_{\mathrm{eff}}\in[0.5,\,0.9]$,
indicating substantial proximity to the GBM limit; however, two equity
assets (\textsc{ABEV3}, \textsc{BBDC4}) yield $\hat{\alpha}\approx2.45$,
placing them near the white-noise end of the reference band and suggesting
lower intraday spectral persistence in their return dynamics. Among
liquid assets, \textsc{BIT\$D} (Bitcoin futures) yields the lowest
$\hat{\alpha}\approx1.94$ (M5) and $\approx1.82$ (H1), the closest
to the GBM floor ($\hat{\varphi}_{\mathrm{eff}}\approx0.89$), indicating
that this market exhibits the least detectable departure from random-walk
spectral dynamics at the timescales analysed---consistent with high
intraday efficiency in a continuously traded, globally accessible
asset. 

The procedure for identifying these exponents is analogous to focusing
an optical instrument: the exponents are the physical observables,
and $(\Delta\tau_{\mathrm{opt}},\,L_{\mathrm{opt}})$ are the settings
required to resolve them without aberration from finite-window or
finite-step artifacts. For liquid assets with continuous price formation,
these parameters reflect intrinsic dynamical scales of $K(\tau)$,
set by the autocorrelation structure of the operator rather than by
the observation frequency --- a property confirmed by their stability
across M5 and H1, which differ by a factor of 12 in physical time.
For assets with intermittent or structurally low liquidity, the same
criteria yield effective parameters that absorb microstructural contributions.
In these cases $\hat{\beta}_{\mathrm{struct}}$ remains well-defined,
but its interpretation requires accounting for the non-uniform sampling
of the underlying continuous process. Notably, the tail exponent $\hat{\alpha}$
estimated for low-liquidity assets tends to exceed that of liquid
ones --- with values reaching $\hat{\alpha}\approx2.5$--$3.7$
--- indicating that structural transitions in the spectral gap become
rarer and less correlated, moving the return-time statistics toward
the white-noise limit. The liquid and low-liquidity regimes are analysed
separately in Section~\ref{sec:regimes}.

The automated criterion recovers the visually identified plateau region
to within one grid step for all assets in the liquid sample (Table~\ref{tab:params}).
Residual discrepancies of $\sim\!1$--$2$ steps occur for assets
where the M5 and H1 curves are systematically offset but individually
flat, reflecting the intrinsic uncertainty in plateau detection when
inter-scale coherence is limited by microstructural differences between
timeframes. In such cases the reported $\Delta\tau_{\mathrm{opt}}$
should be understood as the lower bound of a plateau interval rather
than a sharp optimum, and the sensitivity of the downstream exponents
to this ambiguity is assessed in Supplementary Section~\ref{sec:supp_sensitivity}.

\subsection{Scope and extensions}

The operator-space framework introduced here is not specific to financial
markets. Any multivariate time series admitting a Hankel embedding
can be analysed within the same formalism, and the cross-scale calibration
procedure applies without modification. The extension to other complex
systems like turbulence, neural population dynamics, ecological and
climate time series, where anomalous diffusion and non-trivial spectral
structure coexist~\cite{metzler2000,bouchaud1990}, represents a
natural direction for future investigation.

\section{Conclusion}

We have introduced a framework in which the temporal covariance operator
$K(\tau)$, constructed via Hankel embedding of the return series,
is the fundamental dynamical object of financial markets. The observed
price return is the spectral projection of $K(\tau)$, decomposable
into a structural signal subspace and a Marchenko-Pastur bulk via
Eq.~(\ref{eq:decomp}). The dynamics of $K(\tau)$ are governed by
a generalised Langevin equation with a quasi-block-diagonal memory
kernel, from which three universal exponents emerge. $\bstruct\approx1.0$
characterises near-Brownian diffusion of the structural modes, with
the small excess over unity attributable to weak inter-block coupling
and Poisson shocks. $\bbulk\approx1.3$--$1.7$ characterises anomalous
superdiffusion of the bulk modes, driven by collective Dyson-type
eigenvalue mixing and a divergent memory integral. $\alpha\approx2$
characterises the universal power-law distribution of inter-regime
waiting times, consistent with first-passage theory at a regular critical
point and with the boundary of the CTRW transition~\cite{bouchaud1990,metzler2000}.

Together with the intrinsic scales $\Lopt$ and $\dtopt$, these exponents
are consistent across 14 assets, five asset classes, and two intraday
timescales. Building on the anomalous diffusion programme of Refs.~\cite{morgado2002,oliveira2019}
and extending the static RMT framework of Refs.~\cite{laloux1999,plerou1999,bouchaud2003},
the present work establishes a direct bridge between the physics of
stochastic operators and the phenomenology of financial markets, and
places financial phase transitions in a well-defined universality
class. The rigorous derivation of $\alpha=2$ from the Dyson-Langevin
dynamics of the spectral edge process, and the extension of the framework
to other complex multivariate systems, remain as well-posed problems
for future work.

\bibliographystyle{naturemag}
\bibliography{spectral_dynamics}
 
\end{document}
